% Potenzen, Zehnerpotenzen und Quadratwurzeln – bank/potenzen-wurzeln/ (zusammenbau v0.4, Heft)
\einheitenkopf[t6]{Potenzen, Zehnerpotenzen und Quadratwurzeln}
\setcounter{aufgabe}{44}

% Anlauf (ohne Prüfkennung)
\begin{aufgabe}{Potenzen – Anlauf}
\begin{teile}
\teil $7^4$ – hoch oder mal?\\ \kreuz{hoch} \\ \kreuz{mal}
\teil $2^6$ als Malkette – wie viel? \leerfeld
\teil $(-4)^3$ – wie viel? \leerfeld
\end{teile}
\end{aufgabe}

\begin{aufgabe}{Potenzen – Prüfungsaufgaben}
\begin{teile}
\teil $2^x = 512$ – welche Zahl ist $x$? \hfill (P10 2020 OS) \\ x = \leerfeld
\teil $2^x = 128$ – welche Zahl ist $x$? \hfill (P10 2020 OS) \\ x = \leerfeld
\teil $4^x = 64$ – welche Zahl ist $x$? \hfill (P10 2024 OS) \\ x = \leerfeld
\teil $4^x = 1\,024$ – welche Zahl ist $x$? \hfill (P10 2024 OS) \\ x = \leerfeld
\teil $9^3$, $3^5$, $2^9$ – welche Zahl ist die größte? Unterstreiche sie. \hfill (P10 2019 OS)
\teil $6^4$, $4^5$, $3^6$ – welche Zahl ist die größte? Unterstreiche sie. \hfill (P10 2019 OS)
\teil $5^{-3} \;\square\; 0{,}01$ – $<$, $=$ oder $>$? \hfill (P10 2022 OS)
\teil $2^{-4} \;\square\; 0{,}05$ – $<$, $=$ oder $>$? \hfill (P10 2022 OS)
\end{teile}
\end{aufgabe}

% Anlauf (ohne Prüfkennung)
\begin{aufgabe}{Zehnerpotenzen – Anlauf}
\begin{teile}
\teil $4{,}7 \cdot 10^{8}$ – größer oder kleiner als eins?\\ \kreuz{größer als eins} \\ \kreuz{kleiner als eins}
\teil $10^{4}$ – ausgeschrieben? \leerfeld
\teil $3{,}6 \cdot 10^{-3}$ – als Dezimalzahl? \leerfeld
\end{teile}
\end{aufgabe}

\begin{aufgabe}{Zehnerpotenzen – Prüfungsaufgaben}
\begin{teile}
\teil $6{,}3 \cdot 10^{5}$ – ausgeschrieben? \hfill (P10 2016 OS) \\ \leerfeld
\teil $2{,}07 \cdot 10^{6}$ – ausgeschrieben? \hfill (P10 2016 OS) \\ \leerfeld
\teil $9{,}1 \cdot 10^{9}$ – ausgeschrieben? \hfill (P10 2016 OS) \\ \leerfeld
\teil Das Gehirn eines Menschen hat etwa $8{,}6 \cdot 10^{10}$ Nervenzellen. Welche Zahl passt? Kreuze an. \hfill (P10 2015 OS)\\ \kreuz{$860\,000\,000\,000$} \\ \kreuz{$86\,000\,000\,000$} \\ \kreuz{$8\,600\,000\,000$}
\teil Auf der Erde leben etwa $8{,}1 \cdot 10^{9}$ Menschen. Welche Zahl passt? Kreuze an. \hfill (P10 2015 OS)\\ \kreuz{$81\,000\,000\,000$} \\ \kreuz{$8\,100\,000\,000$} \\ \kreuz{$810\,000\,000$}
\teil Ein Fußballstadion hat etwa $7{,}1 \cdot 10^{4}$ Plätze. Welche Zahl passt? Kreuze an. \hfill (P10 2015 OS)\\ \kreuz{$710\,000$} \\ \kreuz{$71\,000$} \\ \kreuz{$7\,100$}
\end{teile}
\end{aufgabe}

\begin{aufgabe}{Zehnerpotenzen – Prüfungsaufgaben – weiter}
\begin{teile}
\teil $730\,000 = 7{,}3 \cdot 10^{\square}$ – welche Hochzahl? \hfill (P10 2025 OS) \\ \leerfeld
\teil $406\,000 = 4{,}06 \cdot 10^{\square}$ – welche Hochzahl? \hfill (P10 2025 OS) \\ \leerfeld
\teil $3{,}4 \cdot 10^{-4}$ – als Dezimalzahl? \hfill (P10 2019 OS) \\ \leerfeld
\teil $6{,}8 \cdot 10^{-5}$ – als Dezimalzahl? \hfill (P10 2019 OS) \\ \leerfeld
\teil Ein Speicherchip fasst $512\,000\,000\,000$ Byte. Wie viele Byte fassen $100$ solche Chips? Gib die Zahl ausgeschrieben an. \hfill (P10 2015 OS) \\ \leerfeld Byte
\teil Ein Herz schlägt in einem Menschenleben etwa $2\,800\,000\,000$-mal. Wie oft schlagen die Herzen von $100$ Menschen zusammen? Gib die Zahl ausgeschrieben an. \hfill (P10 2015 OS) \\ \leerfeld -mal
\teil Welche Zahl ist die kleinste? Kreuze an. \hfill (P10 2017 OS)\\ \kreuz{$-4 \cdot 10^{-2}$} \\ \kreuz{$-4 \cdot 10^{3}$} \\ \kreuz{$-4 \cdot 10^{2}$} \\ \kreuz{$-4 \cdot 10^{-1}$}
\teil Welche Zahl ist die kleinste? Kreuze an. \hfill (P10 2017 OS)\\ \kreuz{$-7 \cdot 10^{2}$} \\ \kreuz{$-7 \cdot 10^{-3}$} \\ \kreuz{$-7 \cdot 10^{5}$} \\ \kreuz{$-7 \cdot 10^{-1}$}
\end{teile}
\end{aufgabe}

% Anlauf (ohne Prüfkennung)
\begin{aufgabe}{Quadratwurzeln – Anlauf}
\begin{teile}
\teil $\sqrt{121}$ – geht die Wurzel im Kopf auf?\\ \kreuz{ja, Quadratzahl} \\ \kreuz{nein, Näherungswert}
\teil $\sqrt{49}$ – wie viel? Mache die Probe. \leerfeld , Probe: \leerfeld
\teil $\sqrt{(-8)^2}$ – wie viel? \leerfeld
\end{teile}
\end{aufgabe}

\begin{aufgabe}{Quadratwurzeln – Prüfungsaufgaben}
\begin{teile}
\teil Ein quadratischer Teppich hat einen Flächeninhalt von $6{,}25\,\text{m}^2$. Wie lang ist eine Seite? \hfill (P10 2020 OS) \\ \leerfeld[m]
\teil Ein quadratischer Platz hat einen Flächeninhalt von $1\,600\,\text{m}^2$. Wie lang ist eine Seite? \hfill (P10 2020 OS) \\ \leerfeld[m]
\teil Eine quadratische Fliese hat einen Flächeninhalt von $441\,\text{cm}^2$. Wie lang ist eine Seite? \hfill (P10 2020 OS) \\ \leerfeld[cm]
\end{teile}
\end{aufgabe}

\begin{aufgabe}{Quadratwurzeln – Prüfungsaufgaben – weiter}
\begin{teile}
\steil Welche Aussage ist wahr? Kreuze an. \hfill (P10 2015 OS)\\ \kreuz{$\sqrt{(-7)^2} = -7$} \\ \kreuz{$\sqrt{(-7)^2} = 7$} \\ \kreuz{$\sqrt{(-7)^2}$ ist nicht definiert}
\steil Welche Aussage ist wahr? Kreuze an. \hfill (P10 2015 OS)\\ \kreuz{$\sqrt{(-13)^2} = -13$} \\ \kreuz{$\sqrt{(-13)^2} = 13$} \\ \kreuz{$\sqrt{(-13)^2}$ ist nicht definiert}
\teil Welche Aussage ist wahr? Kreuze an. \hfill (P10 2015 OS)\\ \kreuz{$2{,}25 < \frac{9}{4}$} \\ \kreuz{$\frac{7}{3} > \frac{9}{4}$} \\ \kreuz{$\sqrt{5} > \frac{9}{4}$}
\teil Welche Aussage ist wahr? Kreuze an. \hfill (P10 2015 OS)\\ \kreuz{$1{,}75 > \frac{7}{4}$} \\ \kreuz{$\frac{9}{5} < \frac{7}{4}$} \\ \kreuz{$\sqrt{3} < \frac{7}{4}$}
\teil $\sqrt{3}$; $-\frac{1}{4}$; $1{,}7$; $-0{,}26$ – aufsteigend geordnet? \hfill (P10 2014 OS)
\teil $\sqrt{11}$; $-\frac{3}{5}$; $3{,}3$; $-0{,}62$ – aufsteigend geordnet? \hfill (P10 2014 OS)
\teil Eine Leiter ist $7{,}5$ m lang. Ihr Fuß steht $2{,}1$ m von der Hauswand entfernt. Wie hoch reicht sie an der Wand? \hfill (P10 2024 OS) \\ \leerfeld[m]
\teil Eine Rampe ist $4$ m lang und überbrückt waagerecht $3{,}6$ m. Wie groß ist der Höhenunterschied? Runde auf zwei Stellen nach dem Komma. \hfill (P10 2024 OS) \\ \leerfeld[m]
\steil Eine zylinderförmige Dose soll $750\,\text{cm}^3$ fassen und $9$ cm hoch sein. Wie groß ist der Radius? Runde auf zwei Stellen nach dem Komma. \hfill (P10 2022 OS) \\ \leerfeld[cm]
\steil Ein zylinderförmiges Glas fasst $330$ ml und ist $8$ cm hoch. Wie groß ist der Radius? Runde auf zwei Stellen nach dem Komma. \hfill (P10 2022 OS) \\ \leerfeld[cm]
\steil $f(x) = x^2 - 8x + 13$ – Nullstellen mit der p-q-Formel? Runde auf zwei Stellen nach dem Komma. \hfill (P10 2025 OS) \\ x1 = \leerfeld , x2 = \leerfeld
\steil $g(x) = x^2 + 4x - 1$ – Nullstellen mit der p-q-Formel? Runde auf zwei Stellen nach dem Komma. \hfill (P10 2025 OS) \\ x1 = \leerfeld , x2 = \leerfeld
\end{teile}
\end{aufgabe}
